%% Research design section 13 -- the five neighbouring systems, plus whatever
%% P6.4's fresh search turns up.
%%
%% Helix [R1], Vidur [R2], DistServe [R3], LLMServingSim 2.0 [R4],
%% Frontier [R5]. Two lines each, and both are required: what overlaps, and
%% what differs. A related-work entry with only the second is an advertisement.
%%
%% Section 13's prohibition applies with full force here: Helix already does
%% graph-based placement and scheduling, and heteropilot already has a Top-K.
%% The novelty claimed is the exact definition of redundancy, a compression that
%% does not lose the shared bottleneck, and quality management under a bounded
%% evaluation budget.
\section{Related work}
\label{sec:related}

\paragraph{Placement on heterogeneous clusters.}
Helix~\cite{helix} formulates serving over heterogeneous GPUs and network as a max-flow
problem and jointly optimises model placement and request scheduling. The
overlap is the object: it also represents the cluster as a graph, and this work
does not claim otherwise; nor does it contribute a scheduler. The difference is what the representation is used
for. Helix solves a flow problem to place and route; the search here
enumerates complete placements and folds them by exact equivalence so that an
expensive evaluator is invoked once per class, with the shared boundary in the
equivalence signature so that two placements crossing different contended links
are not folded together. Others search directly for a good configuration on
heterogeneous, decentralised or spot
GPUs~\cite{thunderserve,shuntserve,hexgen,melange,spotserve} or for a
parallelisation plan~\cite{alpa,galvatron}; this work contributes the
elimination discipline and the accounting of what was never evaluated.

\paragraph{Prefill--decode disaggregation.}
DistServe~\cite{distserve} separates the two phases onto different resources
and schedules them independently; Splitwise~\cite{splitwise},
TetriInfer~\cite{tetriinfer} and Mooncake~\cite{mooncake} disaggregate as well,
and Llumnix~\cite{llumnix} reschedules requests across engine instances. This
work generates disaggregated candidates and
prices the transfer between the phases as a flow over the graph, which is what
makes the counterexample in Section~\ref{sec:candidates} expressible at
all---two prefill groups sending to one decode group differ only in the uplink
they cross. It does not contribute a disaggregation policy, and the hardware
campaign deploys one only through its own experiment harness, because the
deployment backend of HeteroPilot~\cite{heteropilot} has no router
(Section~\ref{sec:limits}).

\paragraph{Simulation.}
Vidur~\cite{vidur} and LLMServingSim~\cite{llmservingsim,llmservingsim2}
provide simulators for LLM serving; Frontier~\cite{frontier,frontier2} argues
for more comprehensive and accurate simulation, including of disaggregated
architectures. The dependence here is direct---the evaluator is LLMServingSim
with an analytical backend---and so is the limitation: that simulator has no
representation of a link graph, so the evaluator sees a placement's own
bottleneck as a bandwidth but not which resources it shares.
Section~\ref{sec:adapter} reports what could not be expressed rather than
folding it into a bandwidth number.
This work measures search accuracy with prediction accuracy held fixed, which
is a different question from the one a simulator paper asks.
